\section{Discusión técnica}
A partir de los resultados obtenidos en las pruebas de rendimiento y de la experiencia derivada del desarrollo de las herramientas, 
es posible sintetizar diversos hallazgos clave sobre la viabilidad y alcance del sistema.
\medskip

En primer término, se demuestra que \textit{OpenVDB} constituye una alternativa sumamente ventajosa para la representación de tensores 
en entornos de computación gráfica. Su principal fortaleza radica en su naturaleza volumétrica unificada, la cual permite encapsular las 
activaciones completas de un tensor multidimensional en un único archivo físico en disco. Este enfoque optimiza sustancialmente el flujo 
de trabajo al mitigar el cuello de botella en las operaciones de entrada/salida (I/O) y reducir drásticamente el volumen de archivos 
generados frente a las secuencias de imágenes convencionales. Asimismo, su esquema de serialización preserva la precisión numérica de 
punto flotante nativa (\texttt{float32}) de las activaciones, eliminando la degradación introducida por la cuantización.
\medskip

En segundo término, la integración nativa de estos volúmenes dentro de Blender habilita un entorno de inspección visual ágil y altamente 
interactivo. La disponibilidad de grupos de nodos procedurales orientados a la generación, transformación y muestreo de densidades 
volumétricas permite explorar las activaciones en el espacio tridimensional y evaluar modificaciones de visualización en tiempo real. 
Debido a restricciones arquitectónicas de Blender, implementar un esquema análogo fundamentado en secuencias de imágenes 2D demandaría 
una sobrecarga sustancial de configuración manual o el desarrollo de \textit{scripts} auxiliares complejos, lo cual ralentizaría sensiblemente 
la iteración y el análisis.
\medskip

De este modo, la extensión desarrollada consolida a Blender no solo como un motor de renderizado técnico y generación de material didáctico o 
demostrativo, sino como un entorno de exploración analítica e interpretabilidad (XAI) para modelos de aprendizaje profundo.
\medskip

No obstante, la solución actual presenta ciertas limitaciones técnicas que delimitan su alcance:
\begin{itemize}
    \item \textbf{Rendimiento en el \textit{viewport}:} la fluidez de navegación interactiva se degrada proporcionalmente a la cantidad de datos 
    volumétricos densos evaluados de manera simultánea, experimentando caídas considerables en la tasa de cuadros por segundo (\textit{framerate}) 
    cuando se renderiza la totalidad de la red en modos con sombreado activo.
    \item \textbf{Acoplamiento topológico:} si bien el módulo de extracción en PyTorch goza de una alta modularidad que facilita su extensión 
    a otros modelos, el componente de visualización en Blender se diseñó específicamente para la topología de Tiny YOLOv2. Adaptar este 
    subsistema a modelos más diversos demandaría reestructuraciones estructurales de mayor escala.
    \item \textbf{Soporte para arquitecturas no lineales:} la canalización actual asume un flujo computacional estrictamente secuencial. 
    La incorporación de modelos modernos con bifurcaciones, conexiones residuales (\textit{skip connections}) o mecanismos de atención 
    requeriría ampliar los algoritmos de disposición geométrica procedural para representar grafos acíclicos dirigidos (DAG) 
    complejos sin inducir saturación visual ni desbordamiento de memoria gráfica.
\end{itemize}